\section{条件独立与信息结构}
\label{sec:05-Probability/02-Conditional-Probability/04}

两个学生参加同一场考试。你发现一个规律：学生A得分高的时候，学生B得分也倾向于高。两人之间有什么关联？他们一起复习了？互相抄了答案？

在追问因果关系之前，先考虑一个更简单的解释：考试难度本身。如果那天的试卷偏简单，两个人都容易得高分；偏难的话，两个人都容易得低分。也就是说，A和B的成绩之间没有直接影响——它们共享了一个共同的原因。

上一节我们学了独立性：\(P(A\cap B)=P(A)P(B)\)。用这个工具测试一下。假设试卷难度只有两种：简单（概率 \(1/2\)），此时两人各自以 \(0.9\) 的概率及格；困难（概率 \(1/2\)），此时两人各自以 \(0.3\) 的概率及格。让我们计算A和B都及格的总概率。

给定简单卷：\(P(\text{A及格}\mid\text{简单})=0.9\)，\(P(\text{B及格}\mid\text{简单})=0.9\)。假设给定难度后两人表现独立（各自靠自己的临场发挥），那么
\[
P(\text{A,B都及格}\mid\text{简单})=0.9\times0.9=0.81.
\]
给定困难卷：同理，
\[
P(\text{A,B都及格}\mid\text{困难})=0.3\times0.3=0.09.
\]
全概率公式给出无条件的联合概率：
\[
P(\text{A,B都及格})=\frac12\times0.81+\frac12\times0.09=0.45.
\]
而无条件各自及格的边缘概率是：
\[
P(\text{A及格})=\frac12\times0.9+\frac12\times0.3=0.6,
\]
B同理，也是 \(0.6\)。

现在检查无条件独立性：
\[
P(\text{A及格})P(\text{B及格})=0.6\times0.6=0.36\neq0.45.
\]
不独立。但之前我们明明假设了——给定难度后——两人表现独立。无条件和有条件给出了相反的结论。

\subsection{条件独立：在某个层次上分解}

这里的现象是：在每一层"难度"内部，A和B的及格与否可以按乘积分解；但把各层混合起来看，分解就不成立了。我们需要一个概念来精确描述"在给定C的情况下，A和B独立"。

\begin{definition}
给定事件 \(C\) 且 \(P(C)>0\)，事件 \(A\) 与 \(B\) 在条件 \(C\) 下独立，是指
\[
P(A\cap B\mid C)=P(A\mid C)P(B\mid C).
\]
记作 \(A\perp B\mid C\)。
\end{definition}

它说的是：把目光限制在 \(C\) 发生的那些实验结果上，重新归一化概率，在这个缩小的世界里，A和B的联合概率等于各自概率的乘积。

条件独立和无条件独立之间没有必然的包含关系。刚才的例子是"条件独立但不无条件独立"。反过来，"无条件独立但不条件独立"同样可能出现——两种现象会分别在本节后续内容中出现。

对于随机变量，符号 \(X\perp Y\mid Z\) 的含义是：在每个合适的条件层 \(Z=z\) 内，\(X,Y\) 的联合分布都分解为各自条件分布的乘积。这个记号是概率图模型和因果推断的基本语法。

\subsection{共同原因：给定原因后，依赖消失}

考试的例子揭示了共同原因的结构。用一个更日常的场景加固理解。

设 \(R\)：是否下雨，\(S\)：街上是否有人撑伞，\(W\)：道路是否湿。直觉上，\(S\) 和 \(W\) 是相关的——看到有人撑伞，你会猜测外面可能下雨了，从而推测道路也可能是湿的。但从物理机制看，撑伞和道路湿之间没有直接的因果链：是下雨同时导致了这两件事。

赋予具体数值。设下雨概率 \(P(R)=0.3\)。下雨时，撑伞概率 \(P(S\mid R)=0.8\)，道路湿概率 \(P(W\mid R)=0.95\)；不下雨时，\(P(S\mid \lnot R)=0.05\)（偶尔有人遮阳），\(P(W\mid\lnot R)=0.02\)（洒水车）。假设给定 \(R\) 后 \(S\) 和 \(W\) 由独立的局部机制决定，那么
\[
P(S\cap W\mid R)=0.8\times0.95=0.76,\qquad
P(S\cap W\mid\lnot R)=0.05\times0.02=0.001.
\]
全概率展开：
\[
P(S\cap W)=0.3\times0.76+0.7\times0.001=0.2287.
\]
无条件边缘：
\[
P(S)=0.3\times0.8+0.7\times0.05=0.275,
\]
\[
P(W)=0.3\times0.95+0.7\times0.02=0.299.
\]
无条件乘积：\(0.275\times0.299\approx0.0822\neq0.2287\)。\(S\) 和 \(W\) 在总体上明显正相关——联合概率几乎是乘积的三倍。

但给定下雨后：
\[
P(S\mid R)P(W\mid R)=0.8\times0.95=0.76=P(S\cap W\mid R).
\]
给定不下雨后：
\[
P(S\mid\lnot R)P(W\mid\lnot R)=0.05\times0.02=0.001=P(S\cap W\mid\lnot R).
\]
所以在每一层天气条件下，\(S\) 和 \(W\) 的乘积都成立：
\[
S\perp W\mid R.
\]

不观测 \(R\) 时，\(S\) 和 \(W\) 通过 \(R\) 间接关联：观察到有人撑伞 → 提高了对"下雨"的信念 → 提高了对"道路湿"的信念。一旦观察到 \(R\) 本身，猜疑链的中间环节被钉死，上下游的信息不再互相传递——你已经知道原因，看一个结果不会再更新对另一个结果的信念。

这就是朴素Bayes分类器的核心假设：给定类别 \(Y\) 后，各特征 \(X_1,\ldots,X_d\) 条件独立：
\[
P(X_1,\ldots,X_d\mid Y)=\prod_{i=1}^{d}P(X_i\mid Y).
\]
这个假设在现实中很少严格成立——特征之间当然会有超越类别的关联。但它把参数数量从指数级降到了线性级，实际中常常给出可用的近似。

\subsection{共同结果：给定结果后，反而产生依赖}

把箭头反过来，局面就变了。

假设某奖学金录取取决于学术能力 \(X\) 或体育能力 \(Y\)——学生只要在任一维度上足够强就能被录取（"或"的门槛）。在全体申请人中，学术能力和体育能力是独立的：擅长数学不意味着擅长跑步。所以
\[
X\perp Y.
\]
赋予数值。假设两种能力各自以独立概率 \(1/2\) 判定为"强"。录取规则：只要任一能力为强就录取（只有两者都弱才被拒）。录取的概率是
\[
P(\text{录取})=1-P(X\text{弱})P(Y\text{弱})=1-0.5\times0.5=0.75.
\]
被拒绝的学生中，\(X,Y\) 都是弱——此时 \(X\) 和 \(Y\) 的关系是确定的，但不再能讨论"在已录取人群中"的关联了。

现在只看被录取的群体。检查 \(X\) 和 \(Y\) 的条件概率。四个原始组合各有概率 \(1/4\)：
\[
\begin{array}{c|c|c}
(X,Y) & \text{概率} & \text{录取？} \\
\hline
(\text{强,强}) & 0.25 & \text{是} \\
(\text{强,弱}) & 0.25 & \text{是} \\
(\text{弱,强}) & 0.25 & \text{是} \\
(\text{弱,弱}) & 0.25 & \text{否}
\end{array}
\]
在已录取的子群体中（条件概率重新归一化，除以 \(0.75\)）：
\[
\begin{aligned}
P(X\text{强}\mid\text{录取}) &= (0.25+0.25)/0.75 = 2/3,\\
P(Y\text{强}\mid\text{录取}) &= (0.25+0.25)/0.75 = 2/3,\\
P(X\text{强},Y\text{强}\mid\text{录取}) &= 0.25/0.75 = 1/3.
\end{aligned}
\]
检查乘积：\((2/3)\times(2/3)=4/9\approx0.444\)，但 \(1/3\approx0.333\)。不相等。在已录取的人群中，
\[
X\not\perp Y\mid \text{录取}.
\]
原本独立的能力，因为条件于共同的结果，变成了负相关——你如果知道一个学生学术弱，就能推断他大概率体育强（否则他怎么被录取的？）。这就是 explaining away。

这种因为条件于共同结果而产生的虚假关联，在数据分析中有个更具体的名字：选择偏差。当你只分析"已住院"的患者、"已通过初筛"的简历、"已被平台推荐"的内容时，你观察到的关联可能完全不存在于总体中——它是选择本身制造出来的。

\subsection{动手算一个具体例子}

用最简单的数值巩固理解。设 \(X,Y\) 为两枚独立公平硬币的结果，各以 \(1/2\) 概率取 \(0\) 或 \(1\)。令
\[
S=X+Y.
\]
无条件时，\(X,Y\) 独立。联合分布均匀分布在四个有序对上：
\[
P(X=x,Y=y)=\frac14,\qquad x,y\in\{0,1\}.
\]
现在给定 \(S=1\)。此时可能的 \((X,Y)\) 只剩下两种：
\[
(1,0)\quad\text{或}\quad(0,1).
\]
这两个结果在原概率中各占 \(1/4\)，在 \(S=1\) 的条件下各占 \((1/4)/(2/4)=1/2\)。计算边缘条件概率：
\[
P(X=1\mid S=1)=\frac12,\qquad
P(Y=1\mid S=1)=\frac12.
\]
但联合条件概率：
\[
P(X=1,Y=1\mid S=1)=0.
\]
乘积：\(1/2\times1/2=1/4\neq0\)。所以
\[
X\not\perp Y\mid S.
\]
给定总和后，原本独立的两枚硬币变成了完全负相关：一个为 \(1\) 则另一个必为 \(0\)。\(S\) 是 \(X,Y\) 的共同结果——它同时依赖于两者，条件于它就打开了依赖通道。

再算一个反方向的确认：\(X,Y\) 独立硬币，无条件独立。若给定 \(X=1\)，\(Y\) 还独立吗？在 \(X=1\) 的条件下，\(Y\) 仍然是 \(0\) 或 \(1\) 各 \(1/2\)：
\[
P(Y=1\mid X=1)=P(Y=1)=\frac12.
\]
这里 \(X\) 是 \(Y\) 的"原因"吗？不是。\(X\) 压根跟 \(Y\) 的机制无关——知道第一个硬币不会改变第二个硬币的概率。条件于一个原因并不打开新的依赖通道。

\subsection{图结构帮助思考方向}

有向图用节点和箭头组织条件独立关系，三种基本结构覆盖了最常见的依赖模式：
\begin{itemize}
\tightlist
\item 链：\(X\to Z\to Y\)。信息沿箭头方向流动；给定中间节点 \(Z\) 后，\(X,Y\) 条件独立。
\item 叉形：\(X\leftarrow Z\rightarrow Y\)。共同原因结构；给定 \(Z\) 后 \(X,Y\) 独立——这节前半部分的例子。
\item 碰撞：\(X\rightarrow Z\leftarrow Y\)。共同结果结构；不给定 \(Z\) 时 \(X,Y\) 独立（默认阻断），给定 \(Z\) 后反而打开依赖通道——这节后半部分的例子。
\end{itemize}

这三种结构是概率图模型的语法构件。它们是组织分解的语言——告诉你联合分布可以怎样拆成局部的条件概率乘积。但要区分：图告诉你的分解能不能被解释为因果方向，取决于额外的结构因果假设。从数据中的相关性不能唯一确定箭头的指向。

\subsection{错误的条件独立假设会系统性扭曲结论}

条件独立之所以重要，不光是因为它帮助化简计算。每一次你把联合概率写成乘积，就是在做一个关于世界结构的断言。

正确利用时，条件独立把描述 \(d\) 个二进制变量所需的参数从 \(2^d\) 量级降到多项式级。错误利用时，它会制造看似精确、实则系统性偏差的推断。朴素Bayes分类器在特征高度冗余时会过度自信——每个特征携带的重复信息被当作独立的新证据重复计数。遗漏隐藏的共同原因会让人误以为发现了直接的因果关联——看到伞和湿地同时出现就以为撑伞弄湿了路。

每当你写出一个概率乘积——不论是在公式推导中，还是在建模假设里——都值得停下来确认：这个分解来自定义逻辑（如条件概率的链式法则），来自实验设计（随机化保证的独立），还是仅仅因为"这样写起来最方便"？

下一单元\hyperref[ch:039]{``随机变量与分布：把复杂结果压缩成可研究的数值''}把事件推广为数值函数。随机变量让我们不再逐个讨论"大于4"之类的具体事件，而去研究整个数值分布、它的位置和散布程度。
